\section{Gestione del file del database}
Il file di dump di un database è tipicamente un file di grosse dimensioni, il quale ha al suo interno diverse tabelle che riportano le credenziali degli utenti, che sono state criptate con qualche algoritmo di hashing. Condividerne il contenuto tra diversi nodi di un architettura HPC, anche in presenza di un filesystem clusterizzato, un problema abbastanza complesso da risolvere. Un primo metodo, intuitivo, è quello di spezzetare il file e distribuirlo in modo più o meno equo tra i nodi, che per quanto ingenuo può essere applicato in più modi diversi

\subsection{Equa suddivisione}
La suddivisione può avvenire in due modi, all'inizio un master assegna ad ogni nodo una porzione di database da attaccare, oppure dinamicamente un nodo chiede al master una porzione di database ogni volta che ha finito quella corrente. Per rendere il modello omogeneo, si considerino 3 momenti: $T(0)$  il momento in cui la rete è stabilita e tutti i nodi sono appena partiti, $T(W)$ il momento in cui tutti i nodi sono già in grado o hanno appena iniziato ad eseguire il loro primo task, $T(E)$ il momento in cui il programma finisce e il risultato è disponibile. 